For $\epsilon>0$, the question asks for a constant number of summands, independent of the prime $p$, sufficient to express every residue modulo $p$ as a sum of reciprocals of integers at most $p^\epsilon$. Repeated summands are permitted by this formulation; denominators divisible by $p$ are, of course, excluded.

The full question has an affirmative answer. Glibichuk's primary abstract gives, for sufficiently large $p$, the stronger distinct-denominator conclusion with
\[
 C_\epsilon=8\bigl(\lfloor1/\epsilon+1/2\rfloor+1\bigr)^2;
\]
see \emph{Combinational properties of sets of residues modulo a prime and the Erd\H{o}s--Graham problem}, \url{https://www.mathnet.ru/eng/mzm2708}. We do not reconstruct its small-$\epsilon$ sum-product argument here. This attempt proves the result for every $\epsilon>1/2$, including distinct denominators for large primes, and the general counting obstruction. The range $0<\epsilon\le1/2$ remains unresolved within this attempt.
